\section{The actual slow-band formula for the complete source test class}
\label{sec:full-test-slow-band}

This section joins the separately proved causal, negative-support,
and norm-one estimates.  Its formula concerns the original weighted
orthogonal PW projections and actual trace-class cutoff operators.
It is not an assignment of negative translations to a one-sided
quotient.  In particular it records the difference between the
slow-band joint limit and the completed full-strip limit.

\paragraph{Notation and actual source families.}
Fix $0<\alpha<1/2$ and $s\ge0$, and put
\[
 \begin{gathered}
 H=L^2(\mathbb R,w(x)\,dx),\qquad
 w(x)=e^{2\alpha|x|}(1+x^2)^{s/2},\\
 d=\tfrac12-\alpha,\qquad u=\tfrac12+\alpha,
 \qquad b_s=\#\{j\in\mathbb Z_{\ge0}:2j<s-1\}.
 \end{gathered}
 \tag{MA-FT.1}
\]
The sector is either the rational trivial sector, a nonprincipal
primitive Dirichlet character, or a nontrivial finite-order
character over $\mathbb Q(\zeta_3)$ in the actual class-number-one
construction of MA-CGP-001--003.  The latter includes the explicit
non-induced character $\psi_{269,3}$ of MA-CH-003.
All zeros below are those of this sector's own primitive completed
$L$-function, with their full multiplicities $m_\rho$; trivial zeros
cancelled by the archimedean completion are not in the divisor.
Write $\lambda_\rho=\rho-1/2$ and
\[
 \begin{aligned}
 r_\rho^+&=
 \begin{cases}
 m_\rho,&\operatorname{Re}\rho<u,\\
 \min(m_\rho,b_s),&\operatorname{Re}\rho=u,\\
 0,&\operatorname{Re}\rho>u,
 \end{cases}\\
 r_\rho^{\rm st}&=
 \begin{cases}
 m_\rho,&d<\operatorname{Re}\rho<u,\\
 \min(m_\rho,b_s),&\operatorname{Re}\rho\in\{d,u\},\\
 0,&\operatorname{Re}\rho\notin[d,u],
 \end{cases}\\
 \delta_\rho^-&=r_\rho^+-r_\rho^{\rm st},\qquad
 \delta_\rho^+=m_\rho-r_\rho^+,\qquad
 R_\alpha(z)=-2\alpha-\overline z.
 \end{aligned}
 \tag{MA-FT.2}
\]
Thus $\delta^-$ consists of full lengths strictly below $d$ and
the missing length at $d$; $\delta^+$ consists of full lengths
strictly above $u$ and the missing length at $u$.
These are nonnegative integers.  They do not assert the existence
of any exterior zero or deficient boundary block.

Let $J_c\subset H_c=\mathbf1_{[c,\infty)}H$ be the actual PW
theta-image closure at support edge $c\le0$, and let
\[
 P_c=M_{\mathbf1_{[c,\infty)}},\qquad
 Q_c=P_{J_c},\qquad P_c^Z=P_c-Q_c,\qquad
 \tau_c(g)=\operatorname{Tr}_H(C_gP_c^Z).
 \tag{MA-FT.3}
\]
These projections are orthogonal in the single space $H$.
For the rational family $Y=e^{-c}$ is the source PW level.
For a character with conductor $\mathfrak f$,
$Q=|D_K|N\mathfrak f$ and $Y=Q\Lambda=e^{-c}$, where
$\Lambda$ is its primitive archimedean band.
MA-GP-010 and MA-CGP-002 identify these with the actual invariant
finite-$S$ global PW ranges when $S'=S$.  No extra unit-shell
transition is silently erased.  For fixed finite data the
Gram-resolvent realization remains MA-GP-008 or MA-CGP.18, with
both of its normal-equation factors.

For $g\in C_c^\infty(\mathbb R)$ with $g(0)=0$, define
\[
 \begin{aligned}
 G_+(z)&=\int_0^\infty g(t)e^{zt}\,dt,&
 G_-(z)&=\int_{-\infty}^0g(t)e^{zt}\,dt,\\
 G(z)&=G_+(z)+G_-(z),&
 \tau_{{\rm st},s}(g)&=\sum_\rho r_\rho^{\rm st}G(\lambda_\rho).
 \end{aligned}
 \tag{MA-FT.4}
\]
The two half-line integrals are scalars; no nonsmooth split of
$g$ is being inserted into an unproved operator functional.

\paragraph{MA-FT-001 (full-test actual limit, with both deficiencies).}
For the actual operators (MA-FT.3),
\[
 \begin{split}
 L_{\alpha,s}(g)
   &:=\lim_{c\to-\infty}\tau_c(g)\\
   &=\tau_{{\rm st},s}(g)
     +\sum_\rho\delta_\rho^-
       \left\{G_+(\lambda_\rho)
                      +G_-(R_\alpha\lambda_\rho)\right\}.
 \end{split}
 \tag{MA-FT.5}
\]
Every series displayed here and below is absolutely convergent.
For any source-thin cutoff path $X,Y\to\infty$, $Y/X\to0$,
and $\delta X\to0$, using the source profiles of MA-AT-001,
\[
 \lim\operatorname{Tr}_H
       (M_{\rho_{X,\delta}(e^\cdot)}C_gQ_{-\log Y})
                    =-L_{\alpha,s}(g).
 \tag{MA-FT.6}
\]
The same formula holds for the common fixed annular profile
$a_X(t)=A(t/X)-A(Xt)$ of MA-LA-014 and MA-CQ-003--004.

Choose a fixed finite set $S$ containing every prime-power
contribution, including the reciprocal contributions, met by
$h(t)=g(\log t)$, and also every conductor or different prime
in the character sectors.  Let $T^S_{X,Y,X^{-1}}(h)$ denote the actual
repaired local quotient operator with that common annular profile;
in the character sector this is the actual operator of MA-CQ,
with the conductor and different normalizations there.
Put $P_\eta(g)=G(1/2)+G(-1/2)$ in the rational trivial sector
and $P_\eta(g)=0$ in the stated nontrivial sectors.  Then
\[
 \begin{split}
 \Delta_{\alpha,s}(g)
 &:=\lim_{\substack{X,Y\to\infty\\Y/X\to0}}
  \left\{\operatorname{Tr}_H(M_{a_X(e^\cdot)}C_gQ_{-\log Y})
       -\left(\operatorname{Tr}T^S_{X,Y,X^{-1}}(h)
                                     -P_\eta(g)\right)\right\}\\
 &=\sum_\rho\delta_\rho^+G(\lambda_\rho)
    +\sum_\rho\delta_\rho^-
       \left\{G_-(\lambda_\rho)
                         -G_-(R_\alpha\lambda_\rho)\right\}.
 \end{split}
 \tag{MA-FT.7}
\]
At the lower edge $\operatorname{Re}\rho=d$ one has
$R_\alpha\lambda_\rho=\lambda_\rho$, so its contribution
to the second sum in (MA-FT.7) cancels exactly.  The lower-edge
length has not been lost from (MA-FT.5), and is not replaced
by a full multiplicity before that cancellation.

\begin{proof}
\emph{Actual operator inputs before a scalar evaluation.}
MA-GW-001--004 give the actual source range, its positive
semigroup, the retained half-plane and strip lengths, and the
ambient trace-class lift.  MA-AT-001 then proves
$C_gP_c^Z\in\mathcal S_1(H)$ for every smooth compact $g$,
without a support-sign condition.  These statements are in
the specified weighted space, not just in an equivalent
unweighted model at a fixed level.

First suppose $g$ vanishes on a neighbourhood of zero.
Its positive and negative restrictions extend by zero to
smooth compact tests $g_+$ and $g_-$ with strictly positive
and strictly negative support.  The actual causal trace of
MA-SB-005/006 and MA-GW gives
\[
 \tau_c(g_+)=\sum_\rho r_\rho^+G_+(\lambda_\rho).
 \tag{MA-FT.8}
\]
For the negative test, MA-AT-003, including its proved
character corollary (MA-AT.73), gives
\[
 \lim_{c\to-\infty}\tau_c(g_-)
  =\sum_\rho r_\rho^{\rm st}G_-(\lambda_\rho)
      +\sum_\rho\delta_\rho^-G_-(R_\alpha\lambda_\rho).
 \tag{MA-FT.9}
\]
The reflection in this formula comes from the actual
weighted Hilbert adjoint and the trace-norm limit of its
defect.  It does not replace the primitive sector by its
contragredient.  Add the two actual trace identities, use
$r^+=r^{\rm st}+\delta^-$, and obtain (MA-FT.5) in this case.

\emph{Removal of the norm-one neighbourhood.}
Choose $\chi\in C_c^\infty((-2,2))$ equal to one on
$[-1,1]$, with $0\le\chi\le1$, and put
\[
 r_\epsilon(t)=g(t)\chi(t/\epsilon),\qquad
 g_\epsilon=g-r_\epsilon.
 \tag{MA-FT.10}
\]
Each $g_\epsilon$ is an actual smooth compact test vanishing
near zero.  MA-NZ-001 proves the operator estimate
\[
 \eta_g(\epsilon):=
     \sup_{c\le0}\|C_{r_\epsilon}P_c^Z\|_1
                   \longrightarrow0.
 \tag{MA-FT.11}
\]
Thus $|\tau_c(g)-\tau_c(g_\epsilon)|\le\eta_g(\epsilon)$
uniformly in $c$.  For any $\epsilon,\epsilon'>0$, taking
the already established limits for $g_\epsilon$ and
$g_{\epsilon'}$ gives
\[
 |L_{\alpha,s}(g_\epsilon)-L_{\alpha,s}(g_{\epsilon'})|
              \le\eta_g(\epsilon)+\eta_g(\epsilon').
 \tag{MA-FT.12}
\]
These limits are Cauchy as $\epsilon\downarrow0$.
The same uniform estimate shows $\tau_c(g)$ tends to their
limit: choose $\epsilon$ first to make the two approximation
errors small, and then take $c$ sufficiently negative for
the fixed test $g_\epsilon$.  This is the operator argument
of MA-NZ-002, not convergence of a nonsmooth split in the
smooth test-function topology.

We verify that its scalar value is precisely (MA-FT.5).
For $z=\sigma+i\gamma$ in a fixed vertical strip,
the half-line transforms of $r_\epsilon$ satisfy
\[
 \left|\int_{\mathbb R_\pm}
                r_\epsilon(t)e^{zt}\,dt\right|
       \le C\min\{\epsilon^2,(1+\gamma^2)^{-1}\}.
 \tag{MA-FT.13}
\]
Here $r_\epsilon(0)=0$, while its first derivative at zero
need not vanish.  The $O(\epsilon^2)$ bound follows from
$g(t)=O(t)$ on the shrinking support.  Two integrations
by parts give the other bound with that derivative's
boundary term retained, since
$\|r_\epsilon''\|_1=O(1)$ and
$\|r_\epsilon'\|_1=O(\epsilon)$.
The multiplicity-weighted zero count on each unit-height
interval is $O(\log(2+|\gamma|))$ in these fixed sectors.
Summing (MA-FT.13), first over $|\gamma|\le1/\epsilon$
and then over the remaining intervals, gives
$O(\epsilon\log(2/\epsilon))$.
The reflected exponents have the same imaginary parts
and real parts in another fixed bounded interval.
The same bound therefore applies to every correction
sum.  Each is absolutely convergent for a fixed $g$,
again by two integrations by parts and $g(0)=0$.
All scalar expressions for $g_\epsilon$ consequently
tend to those in (MA-FT.5) for $g$.

\emph{The original cutoff operators.}
MA-AT.8 bounds the difference between the trace in
(MA-FT.6) and $-\tau_{-\log Y}(g)$ by
$\omega_g(\log(X-\delta))\to0$.  Its lower-support
condition is eventually satisfied because $Y/X\to0$,
the support of $g$ is fixed, and $\delta X\to0$.
Its diagonal volume term is zero because $g(0)=0$.
Combining it with (MA-FT.5) proves (MA-FT.6).

For the fixed annular profile, let
$a=\min\operatorname{supp}g$ when $g\ne0$.
The image of $C_gP_c^Z$ is supported in $[c+a,\infty)$.
Eventually $Xe^{c+a}>2$, so the lower term $A(Xe^x)$
vanishes on that image.  The upper term $A(e^x/X)$
is one for $x\le\log X$.  The exact tail factorization
MA-AT.4 therefore bounds the remaining trace error
by $\omega_g(\log X)\to0$.
The compact-position identity term again has trace zero.
This proves the annular version of (MA-FT.6).

The near-zero passage can also be made directly in these
cutoff traces.  For any bounded compact multiplier $M_\rho$
with norm at most one, the actual compact-kernel diagonal
formula and $r_\epsilon(0)=0$ give
\[
 \begin{split}
 \operatorname{Tr}(M_\rho C_{r_\epsilon}Q_c)
     &=-\operatorname{Tr}(M_\rho C_{r_\epsilon}P_c^Z),\\
 |\operatorname{Tr}(M_\rho C_{r_\epsilon}Q_c)|
     &\le\eta_g(\epsilon).
 \end{split}
 \tag{MA-FT.14}
\]
The bound is uniform in both source parameters.  Thus
the use of MA-NZ introduces no additional hidden
order-of-limits assumption at the original operator level.

\emph{The paired local trace.}
The actual local cutoff theorem MA-LA-014, or MA-CQ-003--004
in the character sector, applies to the same fixed
annular profile: $XY\to\infty$ and $Y/X\to0$ are its
two scale conditions.  The geometric limit includes
every prime term by the chosen finite $S$.
The already established full smooth explicit formula,
Meyer's Theorems~5.8 and~5.11 in the rational normalization
of MA-TR/MA-EG, and MA-CH-001 with MA-CQ in the stated
character normalization, give
\[
 \lim\left(\operatorname{Tr}T^S_{X,Y,X^{-1}}(h)
                                -P_\eta(g)\right)
           =-\sum_\rho m_\rho G(\lambda_\rho).
 \tag{MA-FT.15}
\]
This use of the explicit formula needs no information
about which side of $d$ or $u$ a zero lies on.
For characters the conductor term and different factors
are those of MA-CQ; the contact conductor term vanishes
on this class because $g(0)=0$, rather than being omitted
from the local construction.
Subtract (MA-FT.15) from (MA-FT.6), insert (MA-FT.5),
and use $m-r^{\rm st}=\delta^++\delta^-$.
This gives (MA-FT.7).  On the lower edge
$\lambda=-\alpha+i\gamma$, whence
$-2\alpha-\overline\lambda=-\alpha+i\gamma=\lambda$,
proving the final cancellation.
\end{proof}

\paragraph{MA-FT-002 (the full-test arithmetic gate is exact).}
For any one of the stated sectors at fixed $\alpha,s$,
the following conditions are equivalent:
\begin{enumerate}
\item The paired actual local and global limits in
(MA-FT.7) agree for every compact smooth $g$ with $g(0)=0$.
\item They agree for every $g\in C_c^\infty((0,\infty))$.
\item $\delta_\rho^+=0$ for every zero $\rho$.
\item $r_\rho^{\rm st}=m_\rho$ for every zero $\rho$.
\end{enumerate}
One may additionally require both polar moments
$G(1/2)=G(-1/2)=0$ in conditions~1 and~2 without
changing this equivalence.  Requiring these conditions
for every $0<\alpha<1/2$, at any fixed $s\ge0$,
is exactly the Riemann hypothesis for the stated sector.

\begin{proof}
Condition~1 implies~2.  For a strictly positive test,
$G_-=0$, and (MA-FT.7) becomes the trace functional of
the nonnegative integer divisor
\[
                   D^+=\sum_\rho\delta_\rho^+[\lambda_\rho].
 \tag{MA-FT.16}
\]
Its support is locally finite in a fixed vertical
strip and its count is bounded by the sector's full
multiplicity count.  MA-JP-004 therefore proves that
this functional vanishes on all such tests if and
only if $D^+$ is empty.  This proves~2 implies~3.
The same lemma gives a witness with both prescribed
polar moments zero whenever the divisor is nonempty:
nontrivial completed zeros in these sectors do not
have $\lambda=\pm1/2$.  This uses the classical
nonvanishing at the boundary points~0 and~1, already
part of the completed-sector setup, and does not
locate any nontrivial zero within the critical strip.

For completeness, the symmetry needed for~3 implies~4
does not require the character to be self-dual.
The unitary completed functional equation relates
$\Lambda_\eta(z)$ to $\Lambda_{\bar\eta}(1-z)$,
while conjugation relates
$\Lambda_{\bar\eta}(z)$ to
$\overline{\Lambda_\eta(\overline z)}$.
Their combination sends each zero in the same sector
to $1-\overline\rho$, preserving multiplicity.
The rational trivial sector has the same symmetry.
Thus any full or missing length below $d$ is reflected
to exactly the same full or missing length above $u$;
at the two edges the same number $b_s$ is used.
Condition~3 hence also gives $\delta^-=0$, which is~4.
Condition~4 makes both sums in (MA-FT.7) vanish, so it
implies~1.  This also proves the equivalence after
adding the two polar-moment conditions.

If all zeros have real part $1/2$, they lie strictly
inside every strip of positive width, so~4 holds for
every such width, independently of their multiplicity.
Conversely, a zero with real part different from$1/2$
is strictly exterior to a sufficiently narrow positive
strip, where~4 fails.  This proves the last assertion.
\end{proof}

\paragraph{Limit order, source locator, and input boundary.}
The actual completed full-strip trace is
$-\tau_{{\rm st},s}(g)$, by MA-WP-003.
Subtracting it from (MA-FT.6) leaves
\[
 -\sum_\rho\delta_\rho^-
       \{G_+(\lambda_\rho)+G_-(R_\alpha\lambda_\rho)\}.
 \tag{MA-FT.17}
\]
This is the exact order-of-limits correction; it is
not silently identified with zero.  Likewise the
fixed local smoothing profile in this theorem is
not an arbitrary path of thinning local profiles.

The source gap remains at Masullo's blank \S12.3,
p.~514, the asserted comparison in Theorem~12.4.1.1,
pp.~514--519, and the omitted global-to-local bridge
used on p.~520 in Theorem~12.5.1.1.  The printed cutoff
is Definition~9.3.3.1, p.~427; the thin-smoothing
conditions on p.~519 alone do not specify $Y/X\to0$.
MA-JP-001--003 already refute a rate-independent
comparison under only those conditions.  Here the
slow-band path is explicitly supplied and its entire
actual operator limit is proved.  Its agreement with
the unaltered full geometry has precisely the arithmetic
condition MA-FT-002.  No whole-plane zero-location
theorem is used in MA-FT-001 or its equivalence proof.
Using the stipulated theorem $A$ to satisfy that
condition over $\mathbb Q$ or in a Dirichlet sector is
an explicit zero-location input, hence circular for
the requested independent-method test.  The specified
non-induced cubic Hecke sector is not covered by $A$;
its Riemann hypothesis is still the condition, and
$M$ is only a sufficient conditional input.

\paragraph{Trace and finite-length placement.}
The trace-class operators in this theorem use the
same analytic trace framework as Grothendieck's
cote~1, scans~079--080 and~111--112, through the checked
predecessor placements ALPOGE\_BRIDGES BR-TR-01/02.
The integers $m_\rho,r_\rho^{\rm st},\delta_\rho^\pm$
are lengths of finite jet modules supported at
$\lambda_\rho$ in the affine spectral line.
For a finite block, additivity of trace is additivity
on the corresponding finite-length
$\mathbb C[z]$-modules and their $K_0$ classes;
the absolute convergence proved above is what permits
their analytic sum.  The $\lambda$-operations and
Fredholm determinant of the retained actual quotient
are already identified in MA-EG-002, with the exact
SGA~6 placement there.  This theorem does not promote
the reflected half-line expression to a new algebraic
endomorphism, construct an arithmetic weight theorem,
or make the positive source projections a faithful
positive action on the full jet space.  Those stronger
claims do not follow from a trace-class limit.
